\paragraph{\trPointedDisc\trFamily。}\label{sub:families-pointed-discs}
换个方向，取带\trBrane\trLabel{}$(L_0^{\#},\ldots,L_d^{\#})$ 的 $(d{+}1)$ \trPointedDisc{}的一个\trFamily{}$\scr{S}\to\scr{R}$。选好一个\trPerturbationDatum，并考虑与之相关联的一个\trParametrization{}的\trModuliSpace{}$\scr{M}_{\scr{S}}(y_0,\ldots,y_d)$。在\trRegular{}点 $(r,u)$ 处的\trTangentSpace{}现在由\eqref{eq:extended-linearized-operator} 的扩张算子 $D_{\scr{S},r,u}$ 的\trKernel{}给出。把该算子的第一个（有限维）分量\trDeformation{}为零，便得到
\begin{equation}\label{eq:family-moduli-dimension}
\begin{split}
\dim_{(r,u)}\bigl(\scr{M}_{\scr{S}}(y_0,\ldots,y_d)\bigr)
&=\dim_r(\scr{R})+i(y_0)-i(y_1)-\cdots-i(y_d),\\
\bigl(\lambda^{\mathrm{top}}T\scr{M}_{\scr{S}}(y_0,\ldots,y_d)\bigr)_{(r,u)}
&\cong\bigl(\lambda^{\mathrm{top}}T\scr{R}\bigr)_r\otimes o(y_0)\otimes o(y_1)^{\vee}\otimes\cdots\otimes o(y_d)^{\vee}.
\end{split}
\end{equation}
设有两个\trFamily{}$\scr{S}_k\to\scr{R}_k$（$k=1,2$），照前文那样带\trLabel{}$(L_0^{\#},\ldots,L_n^{\#},L_{n+m}^{\#},\ldots,L_d^{\#})$ 与 $(L_n^{\#},\ldots,L_{n+m}^{\#})$，并配备了合适的\trPerturbationDatum。\cref{sub:trees-gluing}与\cref{sub:consistent-perturbation-data}中对这类\trFamily{}引入的\trGluing{}过程（后一小节限于双\trVertex{}树的特殊情形）给出一个诱导的\trFamily{}$\scr{S}\to\scr{R}$，其中 $\scr{R}=(-1;0)\times\scr{R}_1\times\scr{R}_2$。对\trInhomogeneous\trPseudoholomorphicMap，相应的\trGluing{}运算则取\trRegular{}点 $(r_1,u_1)\in\scr{M}_{\scr{S}_1}(y_0,\ldots,y_n,\tilde{y},y_{n+m+1},\ldots,y_d)$ 与 $(r_2,u_2)\in\scr{M}_{\scr{S}_2}(\tilde{y},y_{n+1},\ldots,y_{n+m})$，再配上一个充分小的\trGluingParameter{}$\rho\in(-1;0)$，而给出一个元素 $(r,u)\in\scr{M}_{\scr{S}}(y_0,\ldots,y_d)$。而且，这一运算可扩张成空间 $(-1;0)\times\scr{M}_{\scr{S}_1}\times\scr{M}_{\scr{S}_2}$ 与 $\scr{M}_{\scr{S}}$ 之间的一个\trLocalDiffeomorphism{}$g$。注意，由于其中诉诸\trImplicitFunctionTheorem{}的那一步，$r$ 通常并不等于 $(\rho,r_1,r_2)$；不过这两点是接近的（或者更确切地说，当 $\rho\to 0$ 时\trAsymptotic{}地相互贴近），特别地，可以等同
\begin{equation}\label{eq:family-gluing-tangent}
\bigl(\lambda^{\mathrm{top}}T\scr{R}\bigr)_r
\cong\RR\otimes\bigl(\lambda^{\mathrm{top}}T\scr{R}_1\bigr)_{r_1}\otimes\bigl(\lambda^{\mathrm{top}}T\scr{R}_2\bigr)_{r_2}.
\end{equation}
\eqref{eq:orientation-diagram} 的相应推广则是如下交换图：
\begin{equation}\label{eq:family-gluing-diagram}
\begin{tikzcd}[column sep=4.5em,row sep=3.8em]
\bigl(\lambda^{\mathrm{top}}T\scr{M}_{\scr{S}}\bigr)_{(r,u)}
\ar[r,"\cong"]
&
\bigl(\lambda^{\mathrm{top}}T\scr{R}\bigr)_r\otimes o(y_0)\otimes o(y_1)^{\vee}\otimes\cdots\otimes o(y_d)^{\vee}
\\
{\begin{aligned}[t]
\RR\otimes{}&\bigl(\lambda^{\mathrm{top}}T\scr{M}_{\scr{S}_1}\bigr)_{(r_1,u_1)}\\
\otimes{}&\bigl(\lambda^{\mathrm{top}}T\scr{M}_{\scr{S}_2}\bigr)_{(r_2,u_2)}
\end{aligned}}
\ar[u,"(\lambda^{\mathrm{top}}Dg)_{(\rho,r_1,u_1,r_2,u_2)}"']
\ar[r,"\cong"]
&
{\begin{aligned}[t]
&\RR\otimes\bigl(\lambda^{\mathrm{top}}T\scr{R}_1\bigr)_{r_1}\otimes o(y_0)\\
&\otimes\bigotimes_{k\le n}o(y_k)^{\vee}\otimes o(\tilde{y})^{\vee}\otimes\bigotimes_{k>n+m}o(y_k)^{\vee}\\
&\otimes\bigl(\lambda^{\mathrm{top}}T\scr{R}_2\bigr)_{r_2}\otimes o(\tilde{y})\otimes\bigotimes_{n<k\le n+m}o(y_k)^{\vee}
\end{aligned}}
\ar[u]
\end{tikzcd}
\end{equation}
其中右边的 $\uparrow$ 箭头把\trPairing{}$o(\tilde{y})^{\vee}\otimes o(\tilde{y})\cong\RR$ 与\eqref{eq:family-gluing-tangent} 组合起来。此图在相差 $(-1)^{\circ}$ 的意义下交换，其中现在 $\circ=\dim(\scr{R}_2)\cdot\mathrm{index}(D_{\scr{S}_1,u_1})+\bigl(\sum_{k>n+m}i(y_k)\bigr)\cdot\mathrm{index}(D_{\scr{S}_2,u_2})$。注意，这里的\trIndex{}是诸\trSurface{}$S_k=(\scr{S}_k)_{r_k}$ 上寻常\trLinearization{}算子的\trIndex，而不是其涉及参数空间的扩张版本。
